\section{Results}
\label{sec:results}

\subsection{Primary recovery result}

The Gamma-noise model recovered the prescribed transmission path more
accurately than the constant-\(B\) model on the primary outcome.  Mean
task-level RMSE was 0.653 for Gamma-noise and 1.199 for constant-\(B\).
Gamma-noise had lower RMSE in 195 of 200 paired data sets (97.5\%).  The
task-by-task comparison in Figure~\ref{fig:paired-rmse} shows that the advantage
was widespread rather than being driven by a small subset of simulated
epidemics.

Relative to the constant-model mean, the difference corresponds to an
approximately 45.5\% reduction in mean RMSE.  The five reversals of the paired
ordering are visible above the equality line and prevent interpreting the
result as uniform across all accepted realizations.  The scatter also shows
variation in absolute recovery difficulty, as expected from stochastic
epidemic, reporting, and sampled-trajectory variation.

\begin{figure}[htbp]
\centering
\includegraphics[width=0.68\textwidth]{experiment4/04_paired_RMSE_scatter.pdf}
\caption{Paired RMSE calculated from the sampled Gamma trajectory and repeated
constant estimate for the same 200 accepted data sets.  Points below the
diagonal favor Gamma-noise.}
\label{fig:paired-rmse}
\end{figure}

Table~\ref{tab:errors} places the primary result beside the secondary
signed-error summaries.

\begin{table}[htbp]
\centering
\caption{Mean task-level recovery summaries across 200 paired data sets.}
\label{tab:errors}
\begin{tabular}{lrr}
\toprule
Quantity & Gamma-noise & Constant-\(B\) \\
\midrule
RMSE & 0.653 & 1.199 \\
Absolute task-level mean error & 0.229 & 0.606 \\
Signed mean error, \(t<5\) & -0.086 & -0.424 \\
Signed mean error, \(t\geq5\) & 0.008 & 1.576 \\
\bottomrule
\end{tabular}
\end{table}

\subsection{Secondary recovery summaries}

The mean absolute task-level mean error was 0.229 for Gamma-noise and 0.606
for constant-\(B\).  Gamma-noise had the smaller value in 177 of 200 data sets.
The segment-specific errors explain the main weakness of a single static rate.
Before week~5, mean signed error was \(-0.086\) for Gamma-noise and \(-0.424\)
for constant-\(B\).  From week~5 onward, the corresponding errors were 0.008
and 1.576.  The large positive constant-\(B\) error after the change indicates
that one fitted rate remained well above the prescribed value of 2.

The signed errors also show why the primary RMSE and the absolute mean error
answer different questions.  A constant estimate between 2 and 4 can be too
low in the first regime and too high in the second; averaging those signed
errors permits partial cancellation.  RMSE retains both discrepancies, while
the segment summaries reveal their direction.  Gamma-noise reduced both the
post-change overshoot and the overall displacement, although it did not
produce zero mean error in either segment.

Post-fitting likelihoods gave consistent supporting evidence: mean
particle-filter log likelihood was \(-191.08\) for Gamma-noise and \(-209.23\)
for constant-\(B\), and the saved paired difference favored Gamma-noise in 199
of 200 tasks.  This same-data, unpenalized comparison is a descriptive fitting
diagnostic rather than a formal model-selection test.  The error-distribution
and likelihood-difference figures are retained in
Appendix~\ref{app:diagnostics}.

These results summarize one sampled Gamma trajectory per accepted outbreak
under the specified benchmark.  They do not compare filtering means,
uncertainty intervals, repeated trajectory draws, or performance under other
transmission paths.
